\begin{center}
{\fuenteTitulos\fontsize{14}{17}\selectfont\bfseries\underline{TERMINOS Y DEFINICIONES}}
\end{center}
\vspace{6mm}
{\fontsize{12}{15}\selectfont\bfseries
TPI: Trabajo Práctico Integrador\par
SEP: Sistema Electrónico Programable\par
SSE: Subsistema Electrónico\par
SSEP: Subsistema Electrónico Programable\par
SSW: Subsistema de Software (firmware del SSEP)\par
ESW: Elemento de Software (módulo .c/.h)\par
USW: Unidad de Software (función de un ESW)\par
MCU: Microcontrolador\par
ADC: Conversor Analógico-Digital\par
DMA / GPDMA: Acceso Directo a Memoria / controlador DMA de propósito general del LPC1769\par
LLI: Descriptor de lista enlazada del GPDMA (Linked List Item)\par
PWM: Modulación por Ancho de Pulso\par
UART: Transmisor-Receptor Asíncrono Universal\par
GPIO: Entrada/Salida de Propósito General\par
ISR: Rutina de Servicio de Interrupción\par
NVIC: Controlador de Interrupciones Anidadas Vectorizadas del Cortex-M3\par
PI: Regulador Proporcional-Integral\par
RPM: Revoluciones por minuto\par
CMSIS: Interfaz estándar de software para Cortex-M y biblioteca de drivers LPC17xx\par
CMSIS-DAP: Sonda de depuración y programación integrada en la placa OM13085\par
Por mil: unidad del mando del motor; 1000 por mil equivale a 100\,\% de ciclo de trabajo\par}
